\cleardoublepage
\section{Variables and Algebraic Expressions}

Welcome to the world of algebra! In this section, we'll learn how to use letters like $x$ and $y$ to represent numbers. These letters are called \textbf{variables}, and they help us write expressions and solve equations. Algebra might seem like a new language, but once you understand the rules, it can actually make math easier.

A \textbf{variable} is a symbol, like $x$ or $y$, that stands for a number. We don't always know what the number is right away, but we can still work with it. An \textbf{algebraic expression} is a mix of numbers, variables, and math operations such as $+$, $-$, $\times$, or $\div$.

\subsection*{Example: Understanding Variables and Expressions}

Consider the expression $3x + 5$. This means "3 times some number, then add 5." If $x = 2$, then $3x + 5 = 3(2) + 5 = 6 + 5 = 11$. This is how algebra works: we write ideas using letters and numbers, then solve them by thinking logically.

Imagine you earn \$10 per hour. If $h$ is the number of hours you work, then your earnings are written as $10h$. If you work 5 hours, that's $10(5) = 50$. Variables help us model real-life situations.

Think of a variable as a mystery box. What number could fit inside? Try different values to see how expressions change. When writing expressions, use colors, arrows, or circles to match terms and operations.

The word \textbf{"algebra"} comes from the Arabic term \textbf{"al-jabr,"} which means "reunion of broken parts." It appeared in a famous math book written over 1,200 years ago by a scholar named Al-Khwarizmi. Many English math words like "algorithm" also come from his name!


\cleardoublepage
\subsection{Short Question (English)}
If $x = 3$, find the value of $2x + 7$.

\fullpageimage{images/q2_1_1-eng.jpg}

\newpage
\subsection{Short Question (Korean)}
If $x = 3$, find the value of $2x + 7$.

\fullpageimage{images/q2_1_1-kor.jpg}

\newpage
\subsection{Short Question (Answer)}
If $x = 3$, find the value of $2x + 7$.

\textbf{Answer:} $2(3) + 7 = 6 + 7 = \mathbf{13}$

\cleardoublepage
\subsection{Long Question (English)}
Write an expression for "five less than triple a number y" and evaluate it when $y = 4$.

\fullpageimage{images/q2_1_2-eng.jpg}

\newpage
\subsection{Long Question (Korean)}
Write an expression for "five less than triple a number y" and evaluate it when $y = 4$.

\fullpageimage{images/q2_1_2-kor.jpg}

\newpage
\subsection{Long Question (Answer)}
Write an expression for "five less than triple a number y" and evaluate it when $y = 4$.

\textbf{Answer:}\\
Expression: $3y - 5$\\
When $y = 4$: $3(4) - 5 = 12 - 5 = \mathbf{7}$

\cleardoublepage
\subsection{Challenge Question (English)}
The cost of renting a bicycle is \$8 plus \$3 per hour. Write an expression for the total cost if you rent it for $h$ hours. Find the cost if you rent it for 5 hours.

\fullpageimage{images/q2_1_3-eng.jpg}

\newpage
\subsection{Challenge Question (Korean)}
The cost of renting a bicycle is \$8 plus \$3 per hour. Write an expression for the total cost if you rent it for $h$ hours. Find the cost if you rent it for 5 hours.

\fullpageimage{images/q2_1_3-kor.jpg}

\newpage
\subsection{Challenge Question (Answer)}
The cost of renting a bicycle is \$8 plus \$3 per hour. Write an expression for the total cost if you rent it for $h$ hours. Find the cost if you rent it for 5 hours.

\textbf{Answer:}\\
Expression: $8 + 3h$\\
If $h = 5$: $8 + 3(5) = 8 + 15 = \mathbf{23}$ dollars

